\section{Conclusion and Future Work}
\label{sec:conclusion}

This paper presented an empirical evaluation and statistical verification of tabular reinforcement learning agents navigating continuous 2D platformer environments. By mapping continuous physics into an invariant 7-tuple discrete representation and combining it with asymmetric progress shaping, our system overcame credit assignment bottlenecks without deep neural networks. Through a comprehensive 5-tier curriculum benchmark across 60,000 simulated episodes, we demonstrated that curriculum transfer is essential for multi-hazard domains: while scratch training failed completely on the Level~5 gauntlet (0.0\% win rate), curriculum transfer enabled effective policy acquisition.

Crucially, our comparative benchmark between off-policy Tabular Q-Learning and Hybrid Expected Q-SARSA revealed a fundamental theoretical boundary: on static obstacle stages, pure Q-learning converges reliably; however, in dynamic adversary domains (Level~4), pure Q-learning experiences catastrophic maximization bias under exploratory jitter (3.0\% win rate). By incorporating the behavior policy expectation directly into the TD bootstrap target, Hybrid Expected Q-SARSA effectively prices in exploration risks near hazard boundaries, achieving a \textbf{98.0\% win rate} with zero adversary collisions ($t(181.9) = 28.45, p < 0.0001$, Cohen's $d = 4.02$). Two-sample Welch's $t$-tests confirmed that curriculum policy transfer and expectation-based bootstrapping deliver statistically superior returns and significant simulation compute savings.

\subsection{Team Reflections and Collaborative Takeaways}
Developing this project across a four-member engineering team yielded several key pedagogical insights:
\begin{itemize}
    \item \textbf{Decoupled Architectural Interfaces}: Establishing strict abstract base contracts in \texttt{core/interfaces.py} at project inception allowed Member~1 (Physics), Member~2 (Agent), Member~3 (Training), and Member~4 (Evaluation) to implement, test, and refactor their respective subsystems independently without merge friction.
    \item \textbf{Value of Early Empirical Telemetry}: Implementing automated logging pipelines early in the development lifecycle enabled rapid diagnostic feedback, exposing reward hacking vulnerabilities in dense shaping, maximization bias in pure Q-learning, and perceptual aliasing in compound hazards before master production training.
\end{itemize}

\subsection{Future Scope and Research Directions}
While tabular and hybrid reinforcement learning provides transparency and sample efficiency within structured state manifolds, several promising research avenues remain:
\begin{enumerate}
    \item \textbf{Continuous Function Approximation and Tile Coding}: Transitioning to Tile Coding, Radial Basis Functions (RBFs), or Deep Q-Networks (DQN) would enable the agent to navigate continuous adversary velocities, non-stationary terrain, and procedurally generated layouts without manual bin tuning.
    \item \textbf{Mitigating Perceptual Aliasing in Compound Hazards}: Integrating multi-entity lookahead sensors or recurrent state memory would resolve spatial occlusions when tall obstacles directly abut open chasms or high/low ghost patrols.
    \item \textbf{Multi-Objective Reward Tuning}: Investigating Pareto-optimal trade-offs between safety margins, speed-running traversal time, and optional item collection could uncover diverse navigation personalities.
\end{enumerate}
